Este capítulo apresenta a avaliação experimental do HyBolt e a compara com os resultados publicados para o HySAIL \cite{amaral2013hysail,amaral2019hysail} e com outros esquemas de \acrshort{PoR}. São avaliados, nesta ordem, a recuperação dos dados e o \emph{overhead} de recepção dos códigos \emph{fountain}, o custo do \acrshort{MAC} polinomial, o desempenho das fases de Codificação, Decodificação e Auditoria, a memória exigida do usuário, a eficácia da Auditoria e da ocultação da estrutura do código e, por fim, a comparação com os trabalhos de referência.

\section{Ambiente Experimental}
\label{sec:ambiente_experimental}

Os experimentos foram executados em um computador com processador Intel Core i7-11390H (4 núcleos, 8 \emph{threads}, até 3,4 GHz) e 8 GB de memória, sob o subsistema Linux do Windows (WSL2), com Python~3.13, NumPy~2.2 e a versão 2.0.0 da biblioteca \texttt{raptorq}, implementada em Rust. Os provedores foram simulados por três diretórios no sistema de arquivos Linux do próprio computador, de modo que os tempos medidos não incluem latência de rede.

Salvo indicação em contrário, foram utilizados os parâmetros padrão do HyBolt: $\lambda = 127$, \emph{chunks} de até $\lambda \cdot 512 = 65.024$ bytes, símbolos de 10\% do \emph{chunk}, redundância $\varepsilon = 0{,}25$, $l = 40$ polinômios de desafio no esquema limitado e ocultação da estrutura habilitada. Com esses parâmetros, cada \emph{chunk} completo gera 14 pacotes de 6.500 bytes, dos quais 3 são pacotes de reparo. Os arquivos de entrada foram gerados com bytes pseudoaleatórios, e, em todas as execuções, o arquivo decodificado foi comparado byte a byte com o original.

Todas as medidas e figuras deste capítulo são reproduzidas pelo comando \texttt{python scripts/thesis\_benchmarks.py all}, que grava os dados brutos em \texttt{logs/thesis/} e as figuras em \texttt{doc/thesis/img/resultados/}. O tempo de cada etapa da codificação é obtido envolvendo, durante a medição, os métodos correspondentes do HyBolt com um cronômetro, de modo que as etapas são medidas sem alterar o código da ferramenta. Medidas de arquivos de até 1~MiB correspondem à média de três execuções; as demais, a uma única execução.

\section{Recuperação e \emph{Overhead} de Recepção}
\label{sec:resultados_overhead}

O primeiro experimento compara o número de símbolos codificados necessários para reconstruir $K$ símbolos de origem com RaptorQ, com Online Codes e com LT Codes. Para o RaptorQ, utilizou-se a biblioteca \texttt{raptorq}: em cada repetição, os $2K$ pacotes gerados para um bloco são entregues ao decodificador em ordem aleatória até que a decodificação seja concluída. Como o projeto não possui codificador de Online Codes, os dois outros códigos foram simulados por um decodificador de \emph{peeling} simbólico, que acompanha apenas quais símbolos são conhecidos. Os LT Codes utilizam a distribuição sóliton robusta com $c = 0{,}1$ e $\delta = 0{,}5$, e os Online Codes, a distribuição de Maymounkov \cite{maymounkov2002online} com $q = 3$. Como essa distribuição exige grau máximo $F$ menor que o tamanho da mensagem composta, $\varepsilon$ foi escolhido como o menor valor, a partir de 0,01, que satisfaz essa condição: 0,35 para $K = 11$, 0,10 para $K = 100$ e 0,019 para $K = 1.000$. Foram realizadas 2.000, 1.000 e 200 repetições para $K = 11$, 100 e 1.000, respectivamente.

\begin{figure}[htbp]
	\centering
	\includegraphics[width=\textwidth]{resultados/reception_overhead.pdf}
	\caption{Fração das decodificações concluídas em função do número de símbolos recebidos, normalizado por $K$, para RaptorQ, Online Codes e LT Codes.}
	\label{fig:reception_overhead}
\end{figure}

A Figura~\ref{fig:reception_overhead} e a Tabela~\ref{tab:overhead} mostram a diferença entre os códigos. O RaptorQ concluiu a decodificação com exatamente $K$ símbolos em 99,5\% das repetições com $K = 11$, em 99,6\% com $K = 100$ e em todas as repetições com $K = 1.000$, e nenhuma repetição exigiu mais de $K + 1$ símbolos, o que é compatível com as probabilidades de falha inferiores a 1\% com $K$ símbolos e a 0,01\% com $K + 1$ símbolos especificadas pela RFC~6330 \cite{luby2011rfc6330}. Os Online Codes exigiram, em média, 9,6\% de símbolos adicionais para $K = 1.000$, 31\% para $K = 100$ e 74\% para $K = 11$, e os LT Codes, 17,6\%, 33\% e 56\%, respectivamente. O \emph{overhead} desses códigos só se aproxima do valor assintótico $(1 + 3\varepsilon)K$ para $K$ grande, enquanto o do RaptorQ independe de $K$.

\begin{table}[htbp]
	\centering
	\caption{Símbolos adicionais necessários à recuperação, em porcentagem de $K$}
	\label{tab:overhead}
	\begin{tabular}{lrrrr}
		\hline
		\textbf{Código} & \textbf{$K$} & \textbf{Média} & \textbf{Percentil 99} & \textbf{Falha com $K$ símbolos} \\
		\hline
		RaptorQ & 11 & 0,05\% & 0\% & 0,5\% \\
		Online Codes & 11 & 73,8\% & 191\% & 99,6\% \\
		LT Codes & 11 & 56,3\% & 209\% & 97,5\% \\
		RaptorQ & 100 & 0,004\% & 0\% & 0,4\% \\
		Online Codes & 100 & 31,1\% & 56\% & 100\% \\
		LT Codes & 100 & 32,6\% & 70\% & 100\% \\
		RaptorQ & 1.000 & 0\% & 0\% & 0\% \\
		Online Codes & 1.000 & 9,6\% & 15\% & 100\% \\
		LT Codes & 1.000 & 17,6\% & 28\% & 100\% \\
		\hline
	\end{tabular}
\end{table}

Essa diferença é decisiva para o uso de \emph{chunks} pequenos. Com os parâmetros padrão, cada \emph{chunk} do HyBolt tem $K = 11$ símbolos de origem, valor para o qual os Online Codes exigiriam quase o dobro de símbolos para uma recuperação confiável. O HySAIL original evita esse problema codificando o arquivo inteiro, com $n$ entre 120 e 1.000 blocos de mensagem \cite{amaral2019hysail}, o que, por sua vez, impede limitar a Redistribuição a uma parte do arquivo.

O segundo experimento mede a tolerância a perdas dentro de um \emph{chunk}. Para cada valor de $\varepsilon$, um bloco de 71.500 bytes foi codificado com símbolos de 6.500 bytes, que a biblioteca alinha a 6.496 bytes, resultando em $K = 12$, e $\lceil \varepsilon K \rceil$ símbolos de reparo. Em seguida, foram descartados aleatoriamente de 0 a $K + r$ pacotes, com 1.000 repetições por quantidade descartada.

\begin{figure}[htbp]
	\centering
	\includegraphics[width=0.85\textwidth]{resultados/loss_tolerance.pdf}
	\caption{\emph{Chunks} recuperados em função da fração de pacotes perdidos, para três valores de redundância. As linhas verticais indicam a densidade de corrupção tolerada $\alpha = r / (K + r)$.}
	\label{fig:loss_tolerance}
\end{figure}

Como mostra a Figura~\ref{fig:loss_tolerance}, a recuperação foi completa sempre que foram perdidos até $r$ pacotes e nula a partir de $r + 1$, com uma única exceção: com $\varepsilon = 0{,}5$ e exatamente $r = 6$ pacotes perdidos, 2 das 1.000 repetições falharam, ao restarem exatamente $K$ pacotes. O RaptorQ comporta-se, portanto, praticamente como um código MDS (\emph{Maximum Distance Separable}), e a densidade de corrupção tolerada por \emph{chunk} é $\alpha = r/(K + r)$, como assumido na Seção~\ref{sec:redundancia}: 14\%, 20\% e 33\% para $\varepsilon = 0{,}1$, 0,25 e 0,5. No HySAIL, em contraste, o número de blocos de verificação necessário é dado por $(1 - \delta)(1 + k\delta)n$ e a probabilidade de falha por $\delta^k$, com $k = 3$ \cite{amaral2013hysail}, o que exige, segundo a Figura~7 de \cite{amaral2019hysail}, cerca de 720 blocos de verificação para $n = 500$ e probabilidade de falha de $10^{-3}$, um \emph{overhead} de 44\%.

\section{Custo do MAC Polinomial}
\label{sec:resultados_mac}

A Figura~\ref{fig:mac_runtime} compara o tempo da divisão $m(x) \bmod P(x)$ por um polinômio de grau 127 nas duas implementações do HyBolt com os valores reportados na Figura~5.5 da tese que propôs o HySAIL \cite{amaral2013hysail}. O esquema limitado representa o bloco como um vetor de coeficientes e o reduz com uma tabela de redução por byte, e o esquema compacto utiliza a aritmética de inteiros de precisão arbitrária do Python, em que o polinômio é representado pelos bits de um inteiro.

\begin{figure}[htbp]
	\centering
	\includegraphics[width=0.85\textwidth]{resultados/mac_runtime.pdf}
	\caption{Tempo de cálculo de $m(x) \bmod P(x)$, com $\deg P = 127$, em função do tamanho do bloco. Os valores do HySAIL foram extraídos da Figura~5.5 de \cite{amaral2013hysail}.}
	\label{fig:mac_runtime}
\end{figure}

No HySAIL, o tempo cresce de forma quadrática com o tamanho do bloco, de 0,06~s para 1~KiB a 62~s para 40~KiB \cite{amaral2013hysail}. Nas duas implementações do HyBolt, o crescimento é linear: um bloco de 40~KiB é reduzido em 46~ms no esquema limitado e em 5,2~ms no esquema compacto, respectivamente cerca de 1.300 e 12.000 vezes menos que no HySAIL. Parte dessa diferença decorre do equipamento e da versão do Python, mas a mudança de comportamento assintótico decorre do algoritmo: a redução byte a byte e a redução sobre inteiros processam cada bit do bloco um número constante de vezes. Para o pacote de 6.500 bytes do HyBolt, um \acrshort{MAC} custa cerca de 0,8~ms no esquema compacto, de modo que o \acrshort{MAC} deixa de ser um gargalo do protocolo.

\section{Codificação e Decodificação}
\label{sec:resultados_codificacao}

As Figuras~\ref{fig:encode_runtime} e~\ref{fig:decode_runtime} apresentam os tempos de codificação, incluindo a distribuição dos pacotes e a gravação dos metadados, e de decodificação para arquivos de 64~KiB a 64~MiB. O esquema limitado foi avaliado apenas até 1~MiB, pois sua codificação leva 81~s para esse tamanho.

\begin{figure}[htbp]
	\centering
	\includegraphics[width=0.85\textwidth]{resultados/encode_runtime.pdf}
	\caption{Tempo de codificação e distribuição em função do tamanho do arquivo, para os esquemas compacto e limitado e sem prova de recuperabilidade.}
	\label{fig:encode_runtime}
\end{figure}

\begin{figure}[htbp]
	\centering
	\includegraphics[width=0.85\textwidth]{resultados/decode_runtime.pdf}
	\caption{Tempo de decodificação em função do tamanho do arquivo. Nos esquemas com prova, cada pacote transferido é verificado antes de ser entregue ao decodificador.}
	\label{fig:decode_runtime}
\end{figure}

Todos os esquemas escalam linearmente com o tamanho do arquivo, como esperado da divisão em \emph{chunks} independentes. Para 64~MiB, a codificação levou 22,2~s com o esquema compacto e 5,4~s sem prova de recuperabilidade, o que corresponde a 2,9 e 11,8~MiB/s, e a decodificação, 15,2 e 4,9~s. No esquema limitado, a codificação de 1~MiB levou 81,1~s, 178 vezes mais que no esquema compacto, e a decodificação, 1,37~s, contra 0,24~s. Para arquivos pequenos, o esquema compacto inclui ainda o custo fixo de gerar o polinômio primitivo secreto de grau 127, que pesa mais nos arquivos menores.

\begin{figure}[htbp]
	\centering
	\includegraphics[width=\textwidth]{resultados/encode_stages.pdf}
	\caption{Decomposição do tempo de codificação de um arquivo de 1~MiB por etapa. A etapa ``outros'' reúne a leitura do arquivo, a geração das chaves e o cálculo do \emph{hash} do arquivo.}
	\label{fig:encode_stages}
\end{figure}

A Figura~\ref{fig:encode_stages} mostra onde esse tempo é gasto. No esquema limitado, 96\% da codificação corresponde aos $l = 40$ \acrshort{MAC}s por pacote. No esquema compacto, a etiqueta responde por 66\%, a cifragem por 20\% e a gravação nos provedores por 8\%, enquanto a codificação RaptorQ ocupa apenas 1\% do total. O custo do código de apagamento tornou-se, portanto, desprezível em relação às demais etapas, e otimizações futuras devem concentrar-se no cálculo das etiquetas e na cifragem, ambos implementados em Python puro; a cifragem, em particular, combina o fluxo de chave com o pacote byte a byte.

\section{Auditoria}
\label{sec:resultados_auditoria}

A Figura~\ref{fig:audit_runtime} apresenta o tempo de uma rodada de auditoria sobre os três provedores, desafiando todos os pacotes armazenados ou uma amostra de $s = 50$ pacotes por provedor.

\begin{figure}[htbp]
	\centering
	\includegraphics[width=0.85\textwidth]{resultados/audit_runtime.pdf}
	\caption{Tempo de uma rodada de auditoria sobre três provedores, com todos os pacotes ou com $s = 50$ pacotes por provedor.}
	\label{fig:audit_runtime}
\end{figure}

No esquema compacto, a auditoria de todos os pacotes cresce linearmente com o arquivo, até 3,7~s para os 14.462 pacotes de 64~MiB, enquanto a auditoria amostral permanece entre 40 e 70~ms para arquivos a partir de 1~MiB, independentemente do tamanho do arquivo. No esquema limitado, a rodada amostral de 1~MiB levou 1,2~s, próxima da rodada completa, de 2,0~s, o que indica um custo fixo elevado por rodada; nesse esquema, cada rodada carrega os metadados, que ocupam 1,6~MB, e os regrava ao registrar o polinômio revelado antes do envio do desafio.

O tempo de auditoria é útil apenas se a rodada detectar a corrupção. Para avaliar essa capacidade, um arquivo de 4~MiB foi codificado com o esquema compacto, e uma fração $\beta$ dos 325 pacotes de um provedor teve um byte alterado antes de cada rodada, com 100 repetições por configuração. A Figura~\ref{fig:audit_detection} compara a fração de rodadas que detectaram a corrupção com a probabilidade teórica de uma amostra sem reposição conter um pacote corrompido, $1 - \binom{N - c}{s} / \binom{N}{s}$, em que $N$ é o número de pacotes do provedor e $c$ o de pacotes corrompidos, que se aproxima de $1 - (1 - \beta)^s$ quando $s \ll N$.

\begin{figure}[htbp]
	\centering
	\includegraphics[width=0.85\textwidth]{resultados/audit_detection.pdf}
	\caption{Fração das rodadas de auditoria que detectaram corrupção em função do número de pacotes amostrados por provedor, para 0,9\% (3 pacotes) e 4,9\% (16 pacotes) de pacotes corrompidos.}
	\label{fig:audit_detection}
\end{figure}

Os valores medidos acompanham a curva teórica dentro da variação esperada para 100 repetições, o que indica que a resposta agregada é rejeitada sempre que a amostra contém um pacote corrompido. Além disso, nenhuma das auditorias realizadas sobre arquivos íntegros nos experimentos da Seção~\ref{sec:resultados_codificacao} acusou corrupção. Com 4,9\% de pacotes corrompidos, uma amostra de 50 pacotes detectou a corrupção em 95\% das rodadas; com 0,9\%, foram necessários 200 pacotes para atingir o mesmo patamar.

\section{Memória do Usuário e Armazenamento}
\label{sec:resultados_memoria}

A Figura~\ref{fig:metadata_size} apresenta o tamanho dos metadados mantidos pelo usuário em relação ao tamanho do arquivo.

\begin{figure}[htbp]
	\centering
	\includegraphics[width=0.85\textwidth]{resultados/metadata_size.pdf}
	\caption{Metadados mantidos pelo usuário, em porcentagem do tamanho do arquivo. A linha pontilhada indica o valor que o esquema limitado ocuparia com os \acrshort{MAC}s empacotados em bits.}
	\label{fig:metadata_size}
\end{figure}

No esquema limitado, os metadados ocupam 153\% de um arquivo de 1~MiB. Mesmo que os \acrshort{MAC}s fossem armazenados em bits, e não com um byte por coeficiente, ocupariam $1{,}25 \, l\lambda / B \approx 12\%$ do arquivo, valor que não diminui com o tamanho do arquivo, pois a memória cresce com o número de pacotes. No esquema compacto, os metadados convergem para 3,1\% do arquivo e correspondem quase inteiramente à localização, ao nome e ao \emph{hash} SHA-256 de cada pacote; as chaves da prova ocupam 64 bytes. Essa informação de localização é necessária a qualquer esquema que distribua o arquivo entre provedores, e poderia ser reduzida derivando os nomes dos pacotes das chaves de nomeação, como já é feito na ocultação, e substituindo os \emph{hashes} individuais pela verificação das etiquetas.

O armazenamento nos provedores corresponde a 1,42 vez o tamanho do arquivo para $\varepsilon = 0{,}25$. Além do fator nominal $14/11$, contribuem o preenchimento do último símbolo de cada \emph{chunk}, pois $11 \times 6.496$ bytes excedem os 65.024 bytes do \emph{chunk}, e, em menor escala, o cabeçalho RaptorQ de 4 bytes, a etiqueta de 16 bytes e a serialização de cada pacote.

\section{Ocultação da Estrutura e Ataque Dirigido}
\label{sec:resultados_ocultacao}

A Seção~\ref{sec:problemas_substituicao} mostrou que, sem a ocultação, um adversário que corrompe $r + 1 = 4$ pacotes de um mesmo \emph{chunk} torna o arquivo irrecuperável. Com a ocultação, o adversário não distingue os \emph{chunks} e, na melhor estratégia disponível, corrompe pacotes escolhidos uniformemente. A Figura~\ref{fig:targeted_attack} compara as duas situações, estimando por simulação, com 400 repetições por ponto, a probabilidade de algum \emph{chunk} receber mais de $r$ pacotes corrompidos.

\begin{figure}[htbp]
	\centering
	\includegraphics[width=0.85\textwidth]{resultados/targeted_attack.pdf}
	\caption{Probabilidade de perda do arquivo em função da fração de pacotes corrompidos, com ocultação (corrupção uniforme) e sem ocultação (linhas pontilhadas: $r + 1$ pacotes de um único \emph{chunk}), para arquivos de 1~MiB, 64~MiB e 1~GiB.}
	\label{fig:targeted_attack}
\end{figure}

Para um arquivo de 1~GiB, a ocultação eleva a fração de pacotes que o adversário precisa corromper de 0,0017\% para cerca de 0,3\%, a partir da qual a perda do arquivo passa a ser possível, e para 1,6\%, a partir da qual ela se torna mais provável que a recuperação, um aumento de mais de duas ordens de grandeza. Ao mesmo tempo, a figura revela uma limitação dos \emph{chunks} pequenos: embora cada \emph{chunk} tolere a perda de 20\% de seus pacotes, um arquivo com 16.514 \emph{chunks} tem alta probabilidade de conter algum \emph{chunk} com mais de três pacotes corrompidos quando apenas 1,6\% do total é corrompido. A fração crítica diminui com o número de \emph{chunks} e é de cerca de 10\% para 1~MiB e de 3\% para 64~MiB.

Essa fração de corrupção é, entretanto, detectável. Com 1,6\% de pacotes corrompidos, uma auditoria de 50 pacotes detecta a corrupção em cada provedor com probabilidade de 55\%, e em pelo menos um dos três provedores com probabilidade de 91\%, o que permite iniciar a Redistribuição antes que algum \emph{chunk} se torne irrecuperável. A margem pode ser ampliada aumentando $\varepsilon$, aumentando o tamanho dos \emph{chunks}, o que aproxima o número de pacotes corrompidos de cada \emph{chunk} de sua média $\beta (K + r)$, ou aumentando o tamanho da amostra de auditoria.

\section{Comparação com o HySAIL}
\label{sec:comparacao_hysail}

O artigo do HySAIL reporta os tempos de codificação e decodificação de arquivos de 50, 100 e 500~MB, obtidos com uma implementação em Python~2.7 em um processador Intel Core i7 de 2,3~GHz, para valores de $n$ entre 120 e 1.000 blocos de mensagem \cite{amaral2019hysail}. A codificação torna-se mais rápida com o aumento de $n$, e a decodificação, mais lenta. A Figura~\ref{fig:hysail_comparison} e a Tabela~\ref{tab:comparacao_hysail} comparam os menores tempos reportados com os do HyBolt para os mesmos tamanhos de arquivo. Para o HyBolt, são apresentados o tempo da codificação RaptorQ isolada, o tempo total sem prova de recuperabilidade e sem ocultação, comparável ao do HySAIL, e o tempo total com o esquema compacto e a ocultação.

\begin{figure}[htbp]
	\centering
	\includegraphics[width=\textwidth]{resultados/hysail_comparison.pdf}
	\caption{Tempos de codificação e decodificação do HySAIL, segundo as Figuras 5 e 6 de \cite{amaral2019hysail}, e do HyBolt para arquivos de 50, 100 e 500~MB. A decodificação não possui etapa RaptorQ isolada.}
	\label{fig:hysail_comparison}
\end{figure}

\begin{table}[htbp]
	\centering
	\caption{Tempos de codificação e decodificação, em segundos, do HySAIL e do HyBolt}
	\label{tab:comparacao_hysail}
	\begin{tabular}{lrrrrrr}
		\hline
		& \multicolumn{3}{c}{\textbf{Codificação}} & \multicolumn{3}{c}{\textbf{Decodificação}} \\
		\textbf{Configuração} & \textbf{50 MB} & \textbf{100 MB} & \textbf{500 MB} & \textbf{50 MB} & \textbf{100 MB} & \textbf{500 MB} \\
		\hline
		HySAIL, melhor $n$ \cite{amaral2019hysail} & 3 & 5 & 21 & 2 & 4 & 27 \\
		HyBolt, RaptorQ isolado & 0,12 & 0,28 & 1,8 & -- & -- & -- \\
		HyBolt, sem prova & 0,8 & 1,7 & 11,7 & 0,6 & 1,2 & 7,0 \\
		HyBolt, compacto e ocultação & 17,4 & 31,7 & 148,2 & 14,1 & 24,1 & 138,1 \\
		\hline
	\end{tabular}
\end{table}

Sem prova de recuperabilidade, o HyBolt codificou os arquivos de 1,8 a 3,8 vezes mais rápido e os decodificou de 3,3 a 3,9 vezes mais rápido que os melhores casos do HySAIL, e a codificação RaptorQ isolada foi de 12 a 26 vezes mais rápida que a dos Online Codes. Essa comparação deve ser interpretada com cautela: os equipamentos são de gerações diferentes, e a biblioteca \texttt{raptorq} é compilada, enquanto a implementação do HySAIL foi escrita em Python e declarada não otimizada \cite{amaral2019hysail}. Os resultados indicam, porém, que a substituição dos Online Codes pelos RaptorQ Codes não introduziu custo computacional relevante, apesar da maior complexidade da decodificação por eliminação gaussiana, pois cada \emph{chunk} é decodificado com apenas 11 símbolos de origem.

Com o esquema compacto e a ocultação, o HyBolt é de 5 a 7 vezes mais lento que os tempos reportados pelo HySAIL. Os tempos do HySAIL não incluem, aparentemente, o cálculo dos \acrshort{MAC}s: segundo a Figura~5.5 de \cite{amaral2013hysail}, um único \acrshort{MAC} de um bloco de 40~KiB levava 62~s, e um arquivo de 50~MB dividido em 1.000 blocos tem blocos de 50~KB, de modo que os $n \cdot l$ \acrshort{MAC}s não caberiam nos 3~s reportados. No HyBolt, as etiquetas e a cifragem correspondem a 86\% do tempo de codificação (Figura~\ref{fig:encode_stages}), e a comparação entre as linhas ``sem prova'' e ``compacto e ocultação'' quantifica o custo das garantias de integridade e de ocultação acrescentadas.

Quanto às demais propriedades, o HyBolt mantém a Distribuição entre múltiplos provedores, a Auditoria com localização dos pacotes corrompidos e a Redistribuição do HySAIL. A Redistribuição deixa de ser local, mas transfere apenas o \emph{chunk} afetado, de no máximo 65.024 bytes, e o \emph{overhead} de recepção passa de 10\% a 74\%, conforme $K$, para menos de 0,05\% (Tabela~\ref{tab:overhead}).

\section{Comparação com Outros Esquemas de PoR}
\label{sec:comparacao_por}

A Tabela~\ref{tab:comparacao_por} compara qualitativamente o HyBolt com os esquemas de \acrshort{PoR} de Juels e Kaliski \cite{juels2007pors} e de Shacham e Waters \cite{shacham2008compact}, a partir da comparação apresentada na tese do HySAIL \cite{amaral2013hysail}, com o próprio HySAIL e com os dois esquemas implementados neste trabalho. Esses trabalhos não reportam medidas de desempenho diretamente comparáveis às deste capítulo, de modo que a comparação se restringe às propriedades dos esquemas.

\begin{table}[htbp]
	\centering
	\caption{Comparação qualitativa entre esquemas de \acrshort{PoR}. $t$ é o número de pacotes e $l$ o número de desafios do modelo limitado}
	\label{tab:comparacao_por}
	\footnotesize
	\begin{tabular}{@{}p{3.3cm}p{2.1cm}p{2.1cm}p{2.0cm}p{2.0cm}p{2.1cm}@{}}
		\hline
		\textbf{Propriedade} & \textbf{Juels--Kaliski} & \textbf{Shacham--Waters} & \textbf{HySAIL} & \textbf{HyBolt limitado} & \textbf{HyBolt compacto} \\
		\hline
		Múltiplos provedores & não & não & sim & sim & sim \\
		Redistribuição & não & não & local & por \emph{chunk} & por \emph{chunk} \\
		Auditorias ilimitadas & não & sim & sim\textsuperscript{a} & não ($l$) & sim \\
		Memória do usuário & chave & chave & $n \cdot l$ \acrshort{MAC}s & $t \cdot l$ \acrshort{MAC}s & chave e localização \\
		Resposta & sentinelas & $\approx$ 1 bloco & $\lambda$ bits & $\lambda$ bits & $\approx$ 1 pacote \\
		Ocultação da estrutura & sim & -- & não & sim & sim \\
		\emph{Overhead} de recepção & -- & -- & 10\% a 74\%\textsuperscript{b} & $< 0{,}05\%$ & $< 0{,}05\%$ \\
		\hline
	\end{tabular}

	\footnotesize\raggedright
	\textsuperscript{a} No modelo ilimitado de \cite{amaral2019hysail}; a Seção~\ref{sec:auditoria_limitada} discute a reutilização de polinômios revelados.\\
	\textsuperscript{b} Online Codes simulados neste trabalho, para $K$ entre 1.000 e 11 (Tabela~\ref{tab:overhead}).
\end{table}

O esquema compacto do HyBolt combina as propriedades dos dois grupos: assim como o esquema de Shacham e Waters, admite auditorias ilimitadas com uma chave de tamanho constante e uma resposta do tamanho de um pacote, e, assim como o HySAIL, distribui o arquivo entre provedores e redistribui os pacotes corrompidos. Como em Juels e Kaliski, a estrutura do código é ocultada por cifragem, o que impede o ataque dirigido a um \emph{chunk}. Em relação ao esquema de Shacham e Waters, a diferença principal é a instanciação do \acrshort{MAC} linear com a redução polinomial em $GF(2)[x]$ do HySAIL, cuja garantia de extração ainda deve ser formalizada (Seção~\ref{sec:esquema_compacto}). Nenhum dos esquemas implementados oferece verificação pública, limitação já apontada para o HySAIL \cite{amaral2019hysail} e tratada pelas variantes com auditoria delegada \cite{armknecht2021outsourcing,shen2024public}.

\section{Síntese}

Os experimentos confirmam que a substituição dos Online Codes pelos RaptorQ Codes reduz o \emph{overhead} de recepção a praticamente zero, inclusive para os \emph{chunks} pequenos que limitam o custo da Redistribuição, e que o custo do código de apagamento tornou-se desprezível no protocolo. As garantias de integridade, contudo, não são obtidas pela substituição direta: o esquema limitado, que a implementa, é duas ordens de grandeza mais lento na codificação e exige do usuário metadados maiores que o próprio arquivo. O esquema compacto resolve essas limitações, com metadados de 3\% do arquivo, auditorias ilimitadas em tempo constante quando amostradas e detecção de corrupção coerente com a análise teórica, e a ocultação da estrutura eleva em mais de duas ordens de grandeza a corrupção necessária para destruir um arquivo grande. Os principais pontos de melhoria identificados são o desempenho do cálculo das etiquetas e da cifragem, que dominam o tempo de codificação, e a sensibilidade dos arquivos com muitos \emph{chunks} pequenos à corrupção uniforme, que deve orientar a escolha conjunta de $\varepsilon$, do tamanho dos \emph{chunks} e do tamanho da amostra de auditoria.
